% !TeX root = ../../金融经济学_王江_讲义及习题.tex
% 06_期望效用函数 课后习题（由章节文件拆分）

\begin{problem}
\problemno{6.1} 证明如下形式的期望效用函数：
\[
    U(c)=u_{0}(c_{0})+\sum_{\omega\in\Omega}\pi_{\omega}u_{1}(c_{1\omega})
\]
满足独立性公理。
\end{problem}
\begin{solution}
设两个计划 $c,c'$ 在状态 $\omega'$ 的消费路径同为 $(c_0,x)$，并且 $c\succeq c'$。由效用表示，
\[
 U(c)-U(c')=
 \sum_{\omega\ne\omega'}\pi_\omega
 \bigl[u_1(c_{1\omega})-u_1(c'_{1\omega})\bigr]\ge0.
\]
把共同路径换成 $(c_0,y)$，记新计划为 $d,d'$。两边共同的 $u_0(c_0)$ 和 $\pi_{\omega'}u_1(y)$ 仍抵消，故
\[
 U(d)-U(d')=U(c)-U(c')\ge0.
\]
因此 $d\succeq d'$，排序保持不变，独立性公理成立。
\end{solution}

\begin{problem}
\problemno{6.2*} 对于由练习 6.1 给出的效用函数表示的偏好，证明连续性意味着 $u_{0}(\cdot)$ 和 $u_{1}(\cdot)$ 都是连续的。
\end{problem}
\begin{solution}
由喜多厌少，$u_0,u_1$ 严格递增，且 $\sum_\omega\pi_\omega=1$。

考虑未来各状态消费均为 $y$ 的计划，其效用为
\[
 F(x,y)=u_0(x)+u_1(y).
\]
原偏好限制在这些计划上仍连续。下面证明：若 $F(x,y)=f(x)+g(y)$ 表示连续偏好，且 $f,g$ 严格递增，则 $f$ 连续。

单调函数在内点具有左右极限，且只有至多可数个不连续点。取 $g$ 的一个内点连续点 $y$。
\begin{enumerate}[label=(\roman*)]
\item 若 $f$ 在 $x$ 有左跳跃，记 $\delta=f(x)-f(x-)>0$。由 $g$ 在 $y$ 连续且严格递增，可取 $h>0$ 使
\[
 0<g(y)-g(y-h)<\delta.
\]
取 $x_n\uparrow x$，则
\[
 F(x_n,y)\le f(x-)+g(y)<f(x)+g(y-h)=F(x,y-h).
\]
由下等值集闭性，令 $n\to\infty$ 得 $(x,y)\preceq(x,y-h)$，与 $g$ 严格递增矛盾。
\item 若 $f$ 在 $x$ 有右跳跃，记 $\delta=f(x+)-f(x)>0$。取 $h>0$ 使
\[
 0<g(y+h)-g(y)<\delta.
\]
对 $x_n\downarrow x$，有 $F(x_n,y)>F(x,y+h)$。上等值集闭性推出 $(x,y)\succeq(x,y+h)$，同样矛盾。
\end{enumerate}
所以左右跳跃均不存在，$f$ 连续；在定义域端点只需相应的单侧论证。令 $(f,g)=(u_0,u_1)$ 得 $u_0$ 连续；交换二者，得 $u_1$ 连续。

\end{solution}

\begin{problem}
\problemno{6.3**} 如果 $a\succeq b$ 意味着对于 $\forall\,\alpha\in[0,1]$ 都有 $\alpha a+(1-\alpha)b\succeq b$，那么称偏好关系 $\succ$ 是弱凸的（weakly convex）。这与第 2 章中的强凸性有所不同。假设偏好关系是连续的并且有练习 6.1 给出的期望效用函数表示（因而 $u_{0}(\cdot)$ 和 $u_{1}(\cdot)$ 都是连续函数）。证明此时强凸性包含弱凸性。
\end{problem}
\begin{solution}
按本题的强凸性条件，若 $a\succ b$，则对任意 $0<\alpha<1$，有
\[
 \alpha a+(1-\alpha)b\succ b.
\]
要证弱凸性，只须补上 $a\sim b$ 的情形；$\alpha=0,1$ 以及 $a\succ b$ 的情形都已成立。

设 $a\sim b$，$0<\alpha<1$，记 $m=\alpha a+(1-\alpha)b$。反设 $m\prec b$，则
\[
 U(m)<U(a)=U(b).
\]
由于 $U$ 连续，沿着从 $a$ 到 $m$ 的线段，由介值定理可取
\[
 d=(1-t)a+tm,\quad 0<t<1,\qquad
 U(m)<U(d)<U(b).
\]
于是 $b\succ d$。另一方面，$m$ 位于 $d$ 与 $b$ 之间。具体地，令
\[
 \beta=\frac{\alpha}{1-t+t\alpha}\in(0,1),
 \qquad m=\beta d+(1-\beta)b.
\]
由 $b\succ d$ 和强凸性，应有 $m\succ d$，与 $U(m)<U(d)$ 矛盾。因此 $m\succeq b$，弱凸性成立。

\end{solution}

\begin{problem}
\problemno{6.4**} 假设定义在随机支付上的效用函数具有练习 6.1 中的形式，其中 $\Omega$ 是任意的状态空间而 $P$ 是相应的概率测度，并且偏好满足弱凸性。
\begin{enumerate}[label=(\alph*)]
    \item 证明 $u_{0}(\cdot)$ 和 $u_{1}(\cdot)$ 都是拟凹的。
    \item 它们是否必为凹的？若否，给出反例。
\end{enumerate}
\end{problem}
\begin{solution}
\begin{enumerate}[label=(\alph*)]
    \item $u$ 拟凹是指：若 $u(x)\ge u(y)$，则对任意 $\lambda\in[0,1]$ 有 $u(\lambda x+(1-\lambda)y)\ge u(y)$。也就是较高的那个函数值不会被凸组合拉到更低的值之下。只证 $u_{1}$，$u_{0}$ 同样。任取 $x,y$ 与 $\lambda\in[0,1]$，并取定 $c_{0}$ 以及 $1$ 期除某一状态 $\omega$ 外的消费。令 $a$ 在状态 $\omega$ 的消费为 $x$、$b$ 在状态 $\omega$ 的消费为 $y$，其余消费与 $a$ 相同。若 $u_{1}(x)\ge u_{1}(y)$，则 $U(a)\ge U(b)$，弱凸性给出 $U(\lambda a+(1-\lambda)b)\ge U(b)$，即 $u_{1}(\lambda x+(1-\lambda)y)\ge u_{1}(y)$。故 $u_{1}$ 拟凹。
    \item 不必凹。取 $u_{1}\equiv 0$、$u_{0}(c)=c^{3}$（$c\ge 0$），概率任意。此时 $U$ 只依赖于 $c_{0}$ 且对 $c_{0}$ 严格递增，故 $U(a)\ge U(b)$ 意味着 $a$ 的 $0$ 期消费不低于 $b$，其凸组合亦然，弱凸性成立。但 $u_{0}(c)=c^{3}$ 在 $[0,\infty)$ 上是凸的，不是凹的。
\end{enumerate}
\end{solution}

